\documentclass{book}
\usepackage[backend=biber,natbib=true,style=authoryear]{biblatex}
\addbibresource{/home/nguyen/1_NQBH/math/bib.bib}
\usepackage[utf8]{inputenc}
\usepackage{float}
\usepackage{graphicx}
\usepackage[colorlinks=true,linkcolor=blue,urlcolor=red,citecolor=magenta]{hyperref}
\usepackage{amsmath,amssymb,amsthm}
\allowdisplaybreaks
\numberwithin{equation}{section}
\newtheorem{assumption}{Assumption}[section]
\newtheorem{lemma}{Lemma}[section]
\newtheorem{corollary}{Corollary}[section]
\newtheorem{definition}{Definition}[section]
\newtheorem{proposition}{Proposition}[section]
\newtheorem{theorem}{Theorem}[section]
\newtheorem{notation}{Notation}[section]
\newtheorem{remark}{Remark}[section]
\newtheorem{example}{Example}[section]
\newtheorem{ques}{Question}[section]
\newtheorem{problem}{Problem}[section]
\newtheorem{conjecture}{Conjecture}[section]
\usepackage[left=0.5in,right=0.5in,top=1.5cm,bottom=1.5cm]{geometry}
\usepackage{fancyhdr}
\pagestyle{fancy}
\fancyhf{}
\addtolength{\headheight}{0pt}% obsolete
\lhead{\itshape \scriptsize \chaptername~\thechapter}
\rhead{\itshape \scriptsize \nouppercase{\leftmark}} %\nouppercase !
\renewcommand{\chaptermark}[1]{\markboth{#1}{}}
\cfoot{\thepage}

\title{Psychological Manipulation}
\author{Collector: Nguyen Quan Ba Hong}
\date{\today}

\begin{document}
\maketitle
\setcounter{secnumdepth}{6}
\setcounter{tocdepth}{6}
\tableofcontents

%------------------------------------------------------------------------------%

\part{Wikipedia's}

\chapter{\href{https://en.wikipedia.org/wiki/Psychological_manipulation}{Wikipedia/Psychological Manipulation}}

\textsf{Examples of \href{https://en.wikipedia.org/wiki/Television}{televised} manipulation can be found in \href{https://en.wikipedia.org/wiki/News_programs}{news programs} that can potentially influence mass audiences. Pictured is the infamous \href{https://en.wikipedia.org/wiki/Dziennik_Telewizyjny}{\textit{Dziennik}} (Journal) news cast, which attempted to slander \href{https://en.wikipedia.org/wiki/Capitalism}{capitalism} in \href{https://en.wikipedia.org/wiki/Polish_People's_Republic}{then-communist} Poland using \href{https://en.wikipedia.org/wiki/Loaded_language}{emotive and loaded language}.}

\textit{Psychological manipulation} is a type of \href{https://en.wikipedia.org/wiki/Social_influence}{social influence} that aims to change the behavior or \href{https://en.wikipedia.org/wiki/Perception}{perception} of others through indirect, \href{https://en.wikipedia.org/wiki/Deceptive}{deceptive}, or underhanded tactics.[``Definition of 'Manipulate'''. www.merriam-webster.com. Retrieved 2019-02-24.]

By advancing the interests of the manipulator, often at another's expense, such methods could be considered exploitative and devious.

%
\href{https://en.wikipedia.org/wiki/Social_influence}{Social influence} is not necessarily negative.

E.g., people such as friends, family and doctors, can try to \href{https://en.wikipedia.org/wiki/Persuade}{persuade} to change clearly unhelpful habits and behaviors.

Social influence is generally perceived to be harmless when it respects the right of the influenced to accept or reject it, and is not unduly coercive.

Depending on the context and \href{https://en.wikipedia.org/wiki/Motivations}{motivations}, social influence may constitute underhanded manipulation.

\section{Theories on successful manipulation}
According to psychology author \href{https://en.wikipedia.org/wiki/George_K._Simon}{George K. Simon}, successful psychological manipulation primarily involves the manipulator:[Simon, George K (1996). \textit{In Sheep's Clothing: Understanding and Dealing with Manipulative People}. ISBN 978-1-935166-30-6.]
\begin{itemize}
    \item Concealing \href{https://en.wikipedia.org/wiki/Aggression}{aggressive} intentions and behaviors and being affable.
    \item Knowing the psychological \href{https://en.wikipedia.org/wiki/Vulnerabilities}{vulnerabilities} of the victim to determine which tactics are likely to be the most effective.
    \item Having a sufficient level of ruthlessness to have no qualms about causing harm to the victim if necessary.
\end{itemize}
Consequently, the manipulation is likely to be accomplished through covert aggressive means.[Simon, George K (1996). In Sheep's Clothing: Understanding and Dealing with Manipulative People. ISBN 978-1-935166-30-6.]

\subsection{According to Braiker}
Harriet B. Braiker (2004) identified the following ways that manipulators \href{https://en.wikipedia.org/wiki/Abusive_power_and_control}{control} their victims:[Braiker, Harriet B. (2004). \textit{Who's Pulling Your Strings? How to Break The Cycle of Manipulation}. ISBN 978-0-07-144672-3.]
\begin{itemize}
    \item \href{https://en.wikipedia.org/wiki/Positive_reinforcement}{Positive reinforcement}: includes \href{https://en.wikipedia.org/wiki/Praise}{praise}, \href{https://en.wikipedia.org/wiki/Superficial_charm}{superficial charm}, superficial \href{https://en.wikipedia.org/wiki/Sympathy}{sympathy} (\href{https://en.wikipedia.org/wiki/Crocodile_tears}{crocodile tears}), excessive apologizing, money, approval, gifts, attention, facial expressions such as a forced laugh or smile, and public recognition.
    \item \href{https://en.wikipedia.org/wiki/Negative_reinforcement}{Negative reinforcement}: involves removing one from a negative situation as a reward.
    \item \href{https://en.wikipedia.org/wiki/Reinforcement#Intermittent_reinforcement_schedules}{Intermittent or partial reinforcement}: Partial or intermittent negative reinforcement can create an effective \href{https://en.wikipedia.org/wiki/Climate_of_fear}{climate of fear} and doubt.
    
    Partial or intermittent positive reinforcement can encourage the victim to persist.
    \item \href{https://en.wikipedia.org/wiki/Punishment_(psychology)}{Punishment}: includes \href{https://en.wikipedia.org/wiki/Nagging}{nagging}, yelling, the \href{https://en.wikipedia.org/wiki/Silent_treatment}{silent treatment}, \href{https://en.wikipedia.org/wiki/Intimidation}{intimidation}, threats, swearing, \href{https://en.wikipedia.org/wiki/Emotional_blackmail}{emotional blackmail}, \href{https://en.wikipedia.org/wiki/Guilt_trip}{guilt trips}, sulking, crying, and \href{https://en.wikipedia.org/wiki/Playing_the_victim}{playing the victim}.
    \item Traumatic 1-trial learning: using verbal abuse, explosive anger, or other intimidating behavior to establish dominance or superiority; even 1 incident of such behavior can \href{https://en.wikipedia.org/wiki/Classical_conditioning}{condition} or train victims to avoid upsetting, confronting or contradicting the manipulator.
\end{itemize}

\subsection{According to Simon}
Simon identified the following manipulative techniques:[Simon, George K (1996). \textit{In Sheep's Clothing: Understanding and Dealing with Manipulative People}. ISBN 978-1-935166-30-6.]
\begin{itemize}
    \item \textbf{Lying (by commission)}: It is hard to tell if somebody is lying at the time they do it, although often the truth may be apparent later when it is too late.
    
    1 way to minimize the chances of being lied to is to understand that some personality types (particularly \href{https://en.wikipedia.org/wiki/Psychopaths}{psychopaths}) are experts at lying and cheating, doing it frequently, and often in subtle ways.
    \item \href{https://en.wikipedia.org/wiki/Lie#Lying_by_omission}{\textbf{Lying by omission}}: This is a subtle form of lying by withholding a significant amount of the truth.
    
    This technique is also used in \href{https://en.wikipedia.org/wiki/Propaganda}{propaganda}.
    \item \href{https://en.wikipedia.org/wiki/Denial}{\textbf{Denial}}: Manipulator refuses to admit that they have done something wrong.
    \item \href{https://en.wikipedia.org/wiki/Rationalization_(psychology)}{\textbf{Rationalization}}: An excuse made by the manipulator for inappropriate behavior.
    
    Rationalization is closely related to \href{https://en.wikipedia.org/wiki/Spin_(public_relations)}{spin}.
    \item \href{https://en.wikipedia.org/wiki/Minimisation_(psychology)}{\textbf{Minimization}}: This is a type of denial coupled with rationalization.
    
    The manipulator asserts that their behavior is not as harmful or irresponsible as someone else was suggesting.
    \item \textbf{Selective inattention or selective \href{https://en.wikipedia.org/wiki/Attention}{attention}}: Manipulator refuses to pay attention to anything that may distract from their agenda.
    \item \href{https://en.wikipedia.org/wiki/Distraction}{\textbf{Diversion}}: Manipulator not giving a straight answer to a straight question and instead being diversionary, steering the conversation onto another topic.
    \item \href{https://en.wikipedia.org/wiki/Evasion_(ethics)}{\textbf{Evasion}}: Similar to diversion but giving irrelevant, rambling, or vague responses.
    \item \textbf{Covert \href{https://en.wikipedia.org/wiki/Intimidation}{intimidation}}: Manipulator putting the victim onto the defensive by using veiled (subtle, indirect or implied) threats.
    \item \textbf{Guilt trip}: A special kind of intimidation tactic.
    
    A manipulator suggests to the \href{https://en.wikipedia.org/wiki/Conscientious}{conscientious} victim that they do not care enough, are too selfish or have it too easy.
    
    This can result in the victim feeling bad, keeping them in a \href{https://en.wikipedia.org/wiki/Doubt}{self-doubting}, \href{https://en.wikipedia.org/wiki/Anxious}{anxious} and \href{https://en.wikipedia.org/wiki/Submissive}{submissive} position.
    \item \href{https://en.wikipedia.org/wiki/Shaming}{\textbf{Shaming}}: Manipulator uses \href{https://en.wikipedia.org/wiki/Sarcasm}{sarcasm} and put-downs to increase fear and self-doubt in the victim.
    
    Manipulators use this tactic to make others feel unworthy and therefore defer to them.
    
    Manipulators can make one feel ashamed for even daring to challenge them.
    
    It is an effective way to foster a sense of inadequacy in the victim.
    \item \href{https://en.wikipedia.org/wiki/Victim_blaming}{\textbf{Vilifying the victim}}: This tactic is a powerful means of putting the victim on the defensive while simultaneously masking the aggressive intent of the manipulator, while the manipulator falsely accuses the victim as being an abuser in response when the victim stands up for or defends themselves or their position.
    \item \textbf{\href{https://en.wikipedia.org/wiki/Victim_playing}{Playing the victim} role}: Manipulator portrays themselves as a victim of circumstance or of someone else's behavior in order to gain \href{https://en.wikipedia.org/wiki/Pity}{pity}, \href{https://en.wikipedia.org/wiki/Sympathy}{sympathy} or evoke \href{https://en.wikipedia.org/wiki/Compassion}{compassion} and thereby get something from another.
    
    Caring and conscientious people often cannot stand to see anyone suffering and the manipulator often finds it easy to play on sympathy to get cooperation.
    \item \textbf{Playing the servant role}: Cloaking a self-serving agenda in the guise of a service to a more noble cause.
    \item \href{https://en.wikipedia.org/wiki/Seduction}{\textbf{Seduction}}: Manipulator uses \href{https://en.wikipedia.org/wiki/Superficial_charm}{charm}, praise, flattery or overtly supporting others in order to get them to lower their defenses and give their \href{https://en.wikipedia.org/wiki/Trust_(social_sciences)}{trust} and \href{https://en.wikipedia.org/wiki/Loyalty}{loyalty} to the manipulator.
    
    They will also offer help with the intent to gain trust and access to an unsuspecting victim they have charmed.
    \item \textbf{\href{https://en.wikipedia.org/wiki/Psychological_projection}{Projecting} the \href{https://en.wikipedia.org/wiki/Blame}{blame} (blaming others)}: Manipulator \href{https://en.wikipedia.org/wiki/Scapegoats}{scapegoats} in often subtle, hard-to-detect ways.
    
    Often, the manipulator will project their own thinking onto the victim, making the victim look like they have done something wrong.
    
    Manipulators will also claim that the victim is the one who is at fault for believing lies that they were conned into believing, as if the victim forced the manipulator to be deceitful.
    
    All blame, except for the part that is used by the manipulator to accept false guilt, is done in order to make the victim feel guilty about making healthy choices, correct thinking and good behaviors.
    
    It is frequently used as a means of psychological and emotional manipulation and control. Manipulators lie about lying, only to re-manipulate the original, less believable story into a ``more acceptable'' truth that the victim will believe.
    
    Projecting lies as being the truth is another common method of control and manipulation.
    
    Manipulators may falsely accuse the victim of ``deserving to be treated that way''.
    
    They often claim that the victim is crazy or abusive, especially when there is evidence against the manipulator.
    \item \textbf{Feigning \href{https://en.wikipedia.org/wiki/Innocence}{innocence}}: Manipulator tries to suggest that any harm done was unintentional or that they did not do something that they were accused of.
    
    Manipulator may put on a look of surprise or indignation.
    
    This tactic makes the victim question their own \href{https://en.wikipedia.org/wiki/Judgment}{judgment} and possibly their own sanity.
    \item \textbf{Feigning \href{https://en.wikipedia.org/wiki/Confusion}{confusion}}: Manipulator tries to play dumb by pretending they do not know what the victim is talking about or is confused about an important issue brought to their attention.
    
    The manipulator intentionally confuses the victim in order for the victim to doubt their own accuracy of perception, often pointing out key elements that the manipulator intentionally included in case there is room for doubt.
    
    Sometimes manipulators will have used cohorts in advance to help back up their story.
    \item \textbf{Brandishing anger}: Manipulator uses anger to brandish sufficient emotional intensity and \href{https://en.wikipedia.org/wiki/Rage_(emotion)}{rage} to shock the victim into submission.
    
    The manipulator is not actually angry, they just put on an act.
    
    They just want what they want and get ``angry'' when denied.
    
    Controlled anger is often used as a manipulation tactic to avoid confrontation, avoid telling the truth or to further hide intent.
    
    There are often threats used by the manipulator of going to the police, or falsely reporting abuses that the manipulator intentionally contrived to scare or intimidate the victim into submission.
    
    Blackmail and other threats of exposure are other forms of controlled anger and manipulation, especially when the victim refuses initial requests or suggestions by the manipulator.
    
    Anger is also used as a defense so the manipulator can avoid telling truths at inconvenient times or circumstances.
    
    Anger is often used as a tool or defense to ward off inquiries or suspicion.
    
    The victim becomes more focused on the anger instead of the manipulation tactic.
    \item \textbf{Bandwagon effect}: Manipulator comforts the victim into submission by claiming (whether true or false) that many people already have done something, and the victim should as well.
    
    Such manipulation can be seen in \href{https://en.wikipedia.org/wiki/Peer_pressure}{peer pressure} situations, often occurring in scenarios where the manipulator attempts to influence the victim into trying drugs or other substances.
\end{itemize}

\subsection{Vulnerabilities exploited by manipulators}
According to Braiker's self-help book,[Braiker, Harriet B. (2004). \textit{Who's Pulling Your Strings? How to Break The Cycle of Manipulation}. ISBN 978-0-07-144672-3.] manipulators exploit the following vulnerabilities (buttons) that may exist in victims:
\begin{itemize}
    \item the desire to please
    \item addiction to earning the approval and acceptance of others
    \item emotophobia (fear of negative emotion; i.e. a fear of expressing anger, frustration or disapproval)
    \item lack of \href{https://en.wikipedia.org/wiki/Assertiveness}{assertiveness} and ability to say no
    \item blurry sense of identity (with soft \href{https://en.wikipedia.org/wiki/Personal_boundaries}{personal boundaries})
    \item low \href{https://en.wikipedia.org/wiki/Self-sufficiency}{self-reliance}
    \item external \href{https://en.wikipedia.org/wiki/Locus_of_control}{locus of control}
\end{itemize}
According to Simon,[Simon, George K (1996). \textit{In Sheep's Clothing: Understanding and Dealing with Manipulative People}. ISBN 978-1-935166-30-6.] manipulators exploit the following vulnerabilities that may exist in victims: 
\begin{itemize}
    \item \textbf{Naïveté} - victim finds it too hard to accept the idea that some people are cunning, devious and ruthless or is ``in \href{https://en.wikipedia.org/wiki/Denial}{denial}'' if they are being victimized.
    \item \textbf{Over-\href{https://en.wikipedia.org/wiki/Conscientiousness}{conscientiousness}} - victim is too willing to give manipulator the benefit of the doubt and see their side of things in which they blame the victim.
    \item \textbf{Low self-\href{https://en.wikipedia.org/wiki/Confidence}{confidence}} - victim is \href{https://en.wikipedia.org/wiki/Doubt}{self-doubting}, lacking in confidence and \href{https://en.wikipedia.org/wiki/Assertiveness}{assertiveness}, likely to go on the defensive too easily.
    \item \textbf{Over-\href{https://en.wikipedia.org/wiki/Intellectualization}{intellectualization}} - victim tries too hard to understand and believes the manipulator has some understandable reason to be hurtful.
    \item \textbf{Emotional \href{https://en.wikipedia.org/wiki/Codependent}{dependency}} - victim has a \href{https://en.wikipedia.org/wiki/Submissive}{submissive} or dependent personality.
    
    The more emotionally dependent the victim is, the more vulnerable they are to being exploited and manipulated.
\end{itemize}
Manipulators generally take the time to scope out the characteristics and vulnerabilities of their victims.

%
Kantor advises in his book \textit{The Psychopathology of Everyday Life: How Antisocial Personality Disorder Affects All of Us}[Kantor, Martin (2006). \textit{The Psychopathology of Everyday Life: How Antisocial Personality Disorder Affects All of Us.} ISBN 978-0-275-98798-5.] that vulnerability to \href{https://en.wikipedia.org/wiki/Psychopathic}{psychopathic} manipulators involves being too: 
\begin{itemize}
    \item \href{https://en.wikipedia.org/wiki/Dependent_personality_disorder}{\textbf{Dependent}} - dependent people need to be loved and are therefore gullible and liable to say yes to something to which they should say no.
    \item \href{https://en.wikipedia.org/wiki/Maturity_(psychological)}{\textbf{Immature}} - has impaired judgment and so tends to believe exaggerated advertising claims.
    \item \href{https://en.wikipedia.org/wiki/Na%C3%AFve}{\textbf{Naïve}} - cannot believe there are dishonest people in the world, or takes it for granted that if there are any, they will not be allowed to prey on others.
    \item \textbf{Impressionable} - overly seduced by \href{https://en.wikipedia.org/wiki/Superficial_charm}{charmers}.
    \item \href{https://en.wikipedia.org/wiki/Trusting}{\textbf{Trusting}} - people who are honest often assume that everyone else is honest.
    
    They are more likely to commit themselves to people they hardly know without checking credentials, etc., and less likely to question so-called experts.
    \item \href{https://en.wikipedia.org/wiki/Carelessness}{\textbf{Carelessness}} - not giving sufficient amount of thought or attention on harm or errors.
    \item \href{https://en.wikipedia.org/wiki/Loneliness}{\textbf{Lonely}} - lonely people may accept any offer of human contact.
    
    A psychopathic stranger may offer human companionship for a price.
    \item \href{https://en.wikipedia.org/wiki/Narcissistic}{\textbf{Narcissistic}} - narcissists are prone to falling for unmerited flattery.
    \item \href{https://en.wikipedia.org/wiki/Impulsivity}{\textbf{Impulsive}} - make snap decisions.
    \item \href{https://en.wikipedia.org/wiki/Altruistic}{Altruistic} - the opposite of psychopathic: too honest, too fair, too empathetic.
    \item \href{https://en.wikipedia.org/wiki/Altruistic}{\textbf{Altruistic}} - the opposite of psychopathic: too honest, too fair, too empathetic.
    \item \href{https://en.wikipedia.org/wiki/Frugal}{\textbf{Frugal}} - cannot say no to a bargain even if they know the reason it is so cheap.
    \item \href{https://en.wikipedia.org/wiki/Materialistic}{\textbf{Materialistic}} - easy prey for loan sharks or get-rich-quick schemes.
    \item \href{https://en.wikipedia.org/wiki/Greed}{\textbf{Greedy}} - the greedy and dishonest may fall prey to a psychopath who can easily entice them to act in an immoral way.
    \item \href{https://en.wikipedia.org/wiki/Masochistic_personality_disorder}{\textbf{Masochistic}} - lack self-respect and so unconsciously let psychopaths take advantage of them.
    
    They think they deserve it out of a sense of guilt.
    \item \textbf{The \href{https://en.wikipedia.org/wiki/Elderly}{elderly}} - the elderly can become fatigued and less capable of multi-tasking.
    
    When hearing a sales pitch they are less likely to consider that it could be a con.
    
    They are prone to giving money to someone with a hard-luck story.
    
    See \href{https://en.wikipedia.org/wiki/Elder_abuse}{elder abuse}.
\end{itemize}

\section{Motivations of manipulators}
Manipulators can have various possible motivations, including but not limited to:[Braiker, Harriet B. (2004). \textit{Who's Pulling Your Strings? How to Break The Cycle of Manipulation}. ISBN 978-0-07-144672-3.]
\begin{itemize}
    \item the need to advance their own purposes and personal gain at virtually any cost to others
    \item a strong need to attain feelings of \href{https://en.wikipedia.org/wiki/Abusive_power_and_control}{power} and superiority in relationships with others
    \item a want and need to feel in \href{https://en.wikipedia.org/wiki/Abusive_power_and_control}{control}
    \item a desire to gain a feeling of power over others in order to raise their perception of \href{https://en.wikipedia.org/wiki/Self-esteem}{self-esteem}
    \item boredom, or growing tired of their surroundings, seeing it as a game more than hurting others
    \item covert agenda, criminal or otherwise, including financial manipulation (often seen when the elderly or unsuspecting, unprotected wealthy are intentionally targeted for the sole purpose of obtaining a victim's financial assets)
    \item not identifying with underlying emotions, commitment phobia, and subsequent rationalization (offender does not manipulate consciously, but rather tries to convince themselves of the invalidity of their own emotions)
    \item lack of self control over impulsive and anti social behavior thus pre-emptive or reactionary manipulation to maintain image
\end{itemize}

\section{Psychopathy}
Main article: \href{https://en.wikipedia.org/wiki/Psychopathy}{Psychopathy}.

%
Being manipulative appears in Factor 1 of the \href{https://en.wikipedia.org/wiki/Hare_Psychopathy_Checklist}{Hare Psychopathy Checklist} (PCL).[Skeem, J. L.; Polaschek, D. L. L.; Patrick, C. J.; Lilienfeld, S. O. (2011). "Psychopathic Personality: Bridging the Gap Between Scientific Evidence and Public Policy". Psychological Science in the Public Interest. 12 (3): 95--162. doi:10.1177/1529100611426706. PMID 26167886. S2CID 8521465.]

\subsection{In the workplace}
Main article: \href{https://en.wikipedia.org/wiki/Psychopathy_in_the_workplace}{Psychopathy in the workplace}.

%
1 approach to \href{https://en.wikipedia.org/wiki/Management}{management} in general identifies a very fine, almost non-existent dividing line between management and manipulation.[Frank, Prabbal (2007). \textit{People Manipulation: A Positive Approach} (2 ed.). New Delhi: Sterling Publishers Pvt. Ltd (published 2009). pp. 3--7. ISBN 978-81-207-4352-6. Retrieved 2019-11-09.]

%
The workplace psychopath may often rapidly shift between emotions - used to manipulate people or to cause high anxiety.[Faggioni M \& White M. \textit{Organizational Psychopaths - Who Are They and How to Protect Your Organization from Them} (2009)]

%
The authors of the book \href{https://en.wikipedia.org/wiki/Snakes_in_Suits:_When_Psychopaths_Go_to_Work}{Snakes in Suits: When Psychopaths Go to Work} describe a 5-phase model of how a typical workplace psychopath climbs to and maintains \href{https://en.wikipedia.org/wiki/Power_(social_and_political)}{power}.

In phase 3 (manipulation) the psychopath will create a scenario of \href{https://en.wikipedia.org/wiki/Cognitive_distortion}{``psychopathic fiction''} where positive information about themselves and negative \href{https://en.wikipedia.org/wiki/Disinformation}{disinformation} about others will be created, where one's role as a part of a network of \href{https://en.wikipedia.org/wiki/Flying_monkeys_(psychology)}{pawns or patrons} will be utilized and one will be \href{https://en.wikipedia.org/wiki/Operant_conditioning}{groomed} into accepting the psychopath's agenda.[Baibak, P; Hare, R. D. \textit{Snakes in Suits: When Psychopaths Go to Work} (2007).]

\section{Antisocial, borderline \& narcissistic personality disorders}
According to \href{https://en.wikipedia.org/wiki/Otto_Kernberg}{Kernberg}, \href{https://en.wikipedia.org/wiki/Antisocial_Personality_Disorder}{antisocial}, \href{https://en.wikipedia.org/wiki/Borderline_Personality_Disorder}{borderline}, and \href{https://en.wikipedia.org/wiki/Narcissistic_Personality_Disorder}{narcissistic} personality disorders are all organized at a borderline level of personality organization,[Kernberg, O (1975). \textit{Borderline Conditions and Pathological Narcissism}. New York: Jason Aronson. ISBN 978-0-87668-205-0.] and the 3 share some common characterological deficits and overlapping personality traits, with deceitfulness and exceptional manipulative abilities being the most common traits among antisocial and narcissism.

Borderline is emphasized by unintentional and dysfunctional manipulation, but stigma towards borderlines being deceitful still wrongfully persists.[Aguirre, Blaise (2016). ``Borderline Personality Disorder: From Stigma to Compassionate Care''. Stigma and Prejudice. Current Clinical Psychiatry. Humana Press, Cham. pp. 133--143. doi:10.1007/978-3-319-27580-2\_8. ISBN 9783319275789.]

Antisocials, borderlines, and narcissists are often \href{https://en.wikipedia.org/wiki/Pathological_lying}{pathological liars}.[Kernberg, O (1975). \textit{Borderline Conditions and Pathological Narcissism}. New York: Jason Aronson. ISBN 978-0-87668-205-0.]

Other shared traits may include pathological narcissism,[9] consistent irresponsibility, \href{https://en.wikipedia.org/wiki/Machiavellianism_(psychology)}{Machiavellianism}, lack of \href{https://en.wikipedia.org/wiki/Empathy}{empathy},[Baron-Cohen, S (2012). \textit{The Science of Evil: On Empathy and the Origins of Cruelty}. Basic Books. pp. 45--98. ISBN 978-0-465-03142-9.] cruelty, \href{https://en.wikipedia.org/wiki/Meanness}{meanness}, impulsivity, proneness to self-harm and addictions,[Casillas, A.; Clark, L.A.k (October 2002). ``Dependency, impulsivity, and self-harm: traits hypothesized to underlie the association between cluster B personality and substance use disorders''. Journal of Personality Disorders. 16 (5): 424--36. doi:10.1521/pedi.16.5.424.22124. PMID 12489309.] interpersonal exploitation, hostility, anger and rage, vanity, emotional instability, rejection sensitivity, perfectionism, and the use of primitive \href{https://en.wikipedia.org/wiki/Defence_mechanisms}{defence mechanisms} that are pathological and narcissistic.

Common \href{https://en.wikipedia.org/wiki/Narcissistic_defences}{narcissistic defences} include \href{https://en.wikipedia.org/wiki/Splitting_(psychology)}{splitting}, \href{https://en.wikipedia.org/wiki/Denial}{denial}, \href{https://en.wikipedia.org/wiki/Psychological_projection}{projection}, \href{https://en.wikipedia.org/wiki/Projective_identification}{projective identification}, \href{https://en.wikipedia.org/wiki/Idealization_and_devaluation}{primitive idealization and devaluation}, \href{https://en.wikipedia.org/wiki/Cognitive_distortion}{distortion} (including \href{https://en.wikipedia.org/wiki/Exaggeration}{exaggeration}, \href{https://en.wikipedia.org/wiki/Minimisation_(psychology)}{minimization} and \href{https://en.wikipedia.org/wiki/Lie}{lies}), and \href{https://en.wikipedia.org/wiki/Omnipotence}{omnipotence}.[13]

%
Psychologist \href{https://en.wikipedia.org/wiki/Marsha_M._Linehan}{Marsha M. Linehan} has stated that people with borderline personality disorder often exhibit behaviors which are not truly manipulative, but are erroneously interpreted as such.[``On Manipulation with the Borderline Personality''. ToddlerTime Network. Retrieved Dec 28, 2014.]

According to her, these behaviors often appear as unthinking manifestations of intense pain, and are often not deliberate as to be considered truly manipulative.

In the \href{https://en.wikipedia.org/wiki/DSM-V}{DSM-V}, manipulation was removed as a defining characteristic of borderline personality disorder.[Aguirre, Blaise (2016). ``\textit{Borderline Personality Disorder: From Stigma to Compassionate Care}''. Stigma and Prejudice. Current Clinical Psychiatry. Humana Press, Cham. pp. 133--143. doi:10.1007/978-3-319-27580-2\_8. ISBN 9783319275789.]

%
Manipulative behavior is intrinsic to narcissists, who use manipulation to obtain power and \href{https://en.wikipedia.org/wiki/Narcissistic_supply}{narcissistic supply}.

Those with antisocial personalities will manipulate for material items, power, revenge, and a wide variety of other reasons.[\href{https://en.wikipedia.org/wiki/Psychological_manipulation#CITEREFAmerican_Psychiatric_Association2000}{American Psychiatric Association 2000}]

\section{Histrionic personality disorder}
Main article: \href{https://en.wikipedia.org/wiki/Histrionic_personality_disorder}{Histrionic personality disorder}.

%
People with histrionic personality disorder are usually high-functioning, both socially and professionally.

They usually have good \href{https://en.wikipedia.org/wiki/Social_skills}{social skills}, despite tending to use them to manipulate others into making them the center of attention.[``Histrionic Personality Disorder''. The Cleveland Clinic. Archived from the original on 3 October 2011. Retrieved 23 November 2011.]

\section{Machiavellianism}
\href{https://en.wikipedia.org/wiki/Machiavellianism_(psychology)}{Machiavellianism} is a term that some \href{https://en.wikipedia.org/wiki/Social_psychology}{social} and \href{https://en.wikipedia.org/wiki/Personality_psychology}{personality} psychologists use to describe a person's tendency to be unemotional, uninfluenced by conventional morality and more prone to \href{https://en.wikipedia.org/wiki/Deceive}{deceive} and manipulate others.

In the 1960s, Richard Christie and Florence L. Geis developed a test for measuring a person's level of Machiavellianism (sometimes referred to as the \textit{Machiavelli test}).[Christie, R., and F. L. Geis. (1970) ``How devious are you? Take the Machiavelli test to find out.'' Journal of Management in Engineering 15.4: 17.] 

\section{See also}
\begin{itemize}
    \item Advertising
    \item Appeal to emotion
    \item Blackmail
    \item Brainwashing
    \item Bullying
    \item Culture of fear
    \item Coercion
    \item Coercive persuasion
    \item Confidence trick
    \item Crowd manipulation
    \item Covert hypnosis
    \item Covert interrogation
    \item Dark triad
    \item Deception
    \item Demagogy
    \item Discrediting tactic
    \item DISC assessment
    \item Dumbing down
    \item Fear mongering
    \item Gaslighting
    \item Half-truth
    \item Internet manipulation
    \item Isolation to facilitate abuse
    \item List of confidence tricks
    \item List of fallacies
    \item Lying
    \item Master suppression techniques
    \item Media manipulation
    \item Mind control
    \item Mobbing
    \item Psychological abuse
    \item Psychological warfare
    \item Sheeple
    \item Social engineering (political science)
    \item Social engineering (security)
    \item Social influence
    \item Whispering campaign
    \item Project MKUltra
\end{itemize}

\section{Further reading}

\subsection{Books}
\begin{itemize}
    \item Admin, AB. \textit{30 Covert Emotional Manipulation Tactics: How Manipulators Take Control In Personal Relationships} (2014).
    \item Alessandra, Tony. \textit{Non-Manipulative Selling} (1992).
    \item Barber, Brian K. \textit{Intrusive Parenting: How Psychological Control Affects Children and Adolescents} (2001).
    \item Bowman, Robert P.; Cooper, Kathy; Miles, Ron; \& Carr, Tom. \textit{Innovative Strategies for Unlocking Difficult Children: Attention Seekers, Manipulative Students, Apathetic Students, Hostile Students} (1998).
    \item Bursten, Ben. \textit{Manipulator: A Psychoanalytic View} (1973).
    \item Crawford, Craig. \textit{The Politics of Life: 25 Rules for Survival in a Brutal and Manipulative World} (2007).
    \item Forward, Susan. \textit{Emotional Blackmail} (1997).
    \item Klatte, Bill \& Thompson, Kate. \textit{It's So Hard to Love You: Staying Sane When Your Loved One Is Manipulative, Needy, Dishonest, or Addicted} (2007).
    \item McCoy, Dorothy. \textit{The Manipulative Man: Identify His Behavior, Counter the Abuse, Regain Control} (2006).
    \item McMillan, Dina L. \textit{But He Says He Loves Me: How to Avoid Being Trapped in a Manipulative Relationship} (2008).
    \item Sasson, Janet Edgette. \textit{Stop Negotiating With Your Teen: Strategies for Parenting Your Angry, Manipulative, Moody, or Depressed Adolescent} (2002).
    \item Stern, Robin. \textit{The Gaslight Effect: How to Spot and Survive the Hidden Manipulation Others Use to Control Your Life} (2008).
    \item Swihart, Ernest W. Jr. \& Cotter, Patrick. \textit{The Manipulative Child: How to Regain Control and Raise Resilient, Resourceful, and Independent Kids} (1998).
\end{itemize}

\subsection{Academic journals}
\begin{itemize}
    \item Aglietta M, Reberioux A, Babiak P. ``Psychopathic manipulation in organizations: pawns, patrons and patsies'', in Cooke A, Forth A, Newman J, Hare R (Eds), \textit{International Perspectives and Psychopathy}, British Psychological Society, Leicester, pp. 12--17. (1996).
    \item Aglietta, M.; Reberioux, A.; Babiak, P. ``Psychopathic manipulation at work'', in Gacono, C.B. (Ed), \textit{The Clinical and Forensic Assessment of Psychopathy: A Practitioner's Guide}, Erlbaum, Mahwah, NJ, pp. 287--311. (2000).
    \item Bursten, Ben. ``The Manipulative Personality'', \textit{Archives of General Psychiatry}, Vol 26 No 4, 318--321 (1972).
    \item Buss DM, Gomes M, Higgins DS, Lauterback K. ``Tactics of Manipulation'', \textit{Journal of Personality and Social Psychology}, Vol 52 No 6 1219--1279 (1987).
    \item Fischer, Alexander/Christian Illies. ``Modulated Feelings: The Pleasurable-Ends-Model of Manipulation'', \textit{Philosophical Inquiries IV/2}, 25-44 (2018).
    \item Hofer, Paul. ``The Role of Manipulation in the Antisocial Personality'', \textit{International Journal of Offender Therapy and Comparative Criminology}, Vol. 33 No 2, 91--101 (1989).
\end{itemize}

\section{External links}
\begin{itemize}
    \item Sarah Strudwick (Nov 16, 2010) \href{https://www.youtube.com/watch?v=PwWBHRKFYCA}{Dark Souls - Mind Games, Manipulation and Gaslighting}\hfill$\square$
\end{itemize}

%------------------------------------------------------------------------------%

\printbibliography[heading=bibintoc]
\end{document}